\documentclass[../cours_2de.tex]{subfiles}
\begin{document}
\setcounter{chapter}{12}
\chapter{Fonctions affines et comparaisons}
\label{sec:ch13}
\begin{objectif}
Interpréter pente et ordonnée à l'origine ; démontrer les variations et le signe d'une fonction affine ; comparer deux modèles.
\end{objectif}
\section{Une variation additive constante}
Une fonction affine s'écrit $f(x)=mx+p$ sur $\R$ ou sur un domaine précisé. Si $p=0$, elle est linéaire ; si $m=0$, elle est constante.
\begin{propriete}
Pour $a\ne b$, le taux d'accroissement d'une fonction affine est constant :
\[
\frac{f(b)-f(a)}{b-a}=m.
\]
\end{propriete}
\begin{demonstration}
$f(b)-f(a)=(mb+p)-(ma+p)=m(b-a)$.
Diviser par $b-a\ne0$ donne le résultat. Réciproquement, si ce quotient vaut toujours $m$, choisir un point $a$ fixé donne $f(x)=mx+f(a)-ma$ : la fonction est affine.
\end{demonstration}
Le coefficient $m$ mesure la variation de la sortie pour une unité d'entrée ; son unité est une unité de sortie divisée par une unité d'entrée. L'ordonnée à l'origine vaut $f(0)=p$ si $0$ est dans le domaine.
\begin{exemple}
Un réservoir contient $120$ litres puis perd $3$ litres par minute :
$V(t)=120-3t$ pour $0\leq t\leq40$. La pente est $-3$ litres par minute. Au-delà de $40$, le modèle doit changer car un volume ne devient pas négatif.
\end{exemple}
\section{Pourquoi une droite ?}
L'ensemble des points $(x,f(x))$ vérifie $y=mx+p$ ; c'est une droite d'après le chapitre 10. Une fonction constante donne une droite horizontale, une fonction linéaire une droite passant par l'origine.
Pour reconstruire l'expression à partir de deux couples $(a,f(a))$ et $(b,f(b))$, calculer d'abord $m$, puis $p=f(a)-ma$.
\begin{exemple}
Si $f(2)=7$ et $f(5)=1$, alors $m=(1-7)/(5-2)=-2$ et $p=7+4=11$.
Donc $f(x)=-2x+11$. Un troisième point permettrait de vérifier la compatibilité de données avec ce modèle, mais ne prouverait pas que le phénomène réel reste affine partout.
\end{exemple}
\section{Variations : la preuve}
Une fonction est strictement croissante sur un intervalle si, lorsque l'entrée augmente, sa sortie augmente strictement ; la version non stricte admet l'égalité.
\begin{theoreme}
$f(x)=mx+p$ est strictement croissante si $m>0$, strictement décroissante si $m<0$, et constante si $m=0$.
\end{theoreme}
\begin{demonstration}
Pour $a<b$, $b-a>0$ et $f(b)-f(a)=m(b-a)$.
Cette différence a donc le signe de $m$. Si elle est positive, $f(a)<f(b)$ ; si elle est négative, l'ordre est inversé ; si elle est nulle, les images sont égales.
\end{demonstration}
Cette démonstration explique le rôle de la pente. L'ordonnée à l'origine déplace la droite verticalement sans changer la variation.
\section{Signe et zéro}
Si $m\ne0$, le seul zéro est $r=-p/m$, et $f(x)=m(x-r)$.
\[
\begin{array}{c|ccc}
x&x<r&r&x>r\\\hline
m>0&-&0&+\\
m<0&+&0&-
\end{array}
\]
Ce tableau découle du signe de $x-r$ et de la règle des signes.
Si $m=0$, le signe est simplement celui de $p$ sur tout le domaine.
\begin{attention}
« Croissante » n'est pas synonyme de « positive ». La fonction $x\mapsto x-10$ est croissante sur $\R$, mais négative pour $x<10$. Le signe compare $f(x)$ à $0$ ; les variations comparent $f(a)$ à $f(b)$.
\end{attention}
\section{Comparer deux fonctions affines}
\begin{exemple}
Deux tarifs sont $A(x)=25+4x$ et $B(x)=10+7x$, pour $x\geq0$.
Leur différence est $A(x)-B(x)=15-3x$.
Ils sont égaux pour $x=5$ ; $A>B$ pour $0\leq x<5$ ; $A<B$ pour $x>5$.
Graphiquement, la droite de $A$ est au-dessus de celle de $B$ avant l'intersection, puis en dessous.
\end{exemple}
\begin{center}
\begin{tikzpicture}
\begin{axis}[width=10cm,height=5.7cm,axis lines=left,xmin=0,xmax=8,ymin=0,ymax=70,xlabel={$x$},ylabel={prix},legend pos=north west]
\addplot[accent,thick,domain=0:8]{25+4*x};\addlegendentry{$A$}
\addplot[black,dashed,domain=0:8]{10+7*x};\addlegendentry{$B$}
\addplot[only marks]coordinates{(5,45)};
\end{axis}
\end{tikzpicture}
\end{center}
\section{Interpolation et prudence}
Entre deux mesures $(a,A)$ et $(b,B)$, l'interpolation affine attribue à $x$ la valeur
$A+(B-A)(x-a)/(b-a)$.
Elle est obtenue en imposant la pente $(B-A)/(b-a)$ et le passage par le premier point.
C'est un choix de modèle : entre deux relevés de température, la température réelle n'évolue pas forcément de façon affine.
\utilite{Les fonctions affines décrivent des accroissements constants. Le taux d'accroissement préparera la dérivée, mais aucun calcul différentiel n'est nécessaire pour leurs preuves.}
\begin{retenir}
La pente explique les variations ; le zéro organise le signe ; une différence permet de comparer deux modèles.
\end{retenir}
\end{document}
